\section{Logarithmic orbit integrals and invariant sections}
\label{g:orbits}\label{rb:torus}

The invariant problem uses the natural action on $-K_V=\det T_V$.
Twisting this action by a character changes the invariant subspace and
is not part of the statement.
Both proofs below construct the two weights required by
Lemma~\ref{g:transfer}: a pushforward-volume weight and a fiber-supremum
weight. The first obtains their plurisubharmonicity from invariant
adjoint rank one; the second obtains it by logarithmic orbit integration.
The second calculation is independent of invariant direct-image
curvature, while its final section construction still uses the
allowed-pole induction already proved.

\begin{theorem}\label{rb:invariant}
Let $V$ be a smooth connected projective variety over $\mathbb C$ with
$H^0(V,\Omega_V^j)=0$ for all $j>0$.  Suppose $-K_V$ has a smooth
Hermitian metric of semipositive curvature.  For any algebraic torus
$T$ acting on $V$, there is an integer $m>0$ such that
\[
                         H^0(V,-mK_V)^T\ne0
\]
for the natural anticanonical action.
\end{theorem}

\subsection{The rank-one quotient construction}

\begin{proof}
Average local weights, with their line-bundle transition rule, over the
compact form of $T$ using the natural linearization.  This gives a
smooth semipositive metric invariant under that compact group.
If the action is trivial, Theorem~\ref{rb:allowed} applies with
$D=0$ and $A=L=-K_V$.  Otherwise take the rational quotient with field
$\mathbb C(V)^T$ \cite{Rosenlicht1956,BellGhiocaReichstein2017}, and its connected-fibre smooth
projective model $Y$.  The invariant field is relatively algebraically
closed in $\mathbb C(V)$: an element algebraic over it has a finite
orbit, and the connected torus acts trivially on that orbit.  General
fibres contain a dense orbit.

Choose the model diagram equivariantly, by modifying the trivially
acted-on base, normalizing the dominating graph, and taking an
equivariant resolution.  Write $p:W\to V$ and $f:W\to Y$ for the smooth resolution maps,
$J=K_{W/V}$, and $\bar V$ for the normal equidimensional intermediate
model of Section~\ref{g:models}. Let $r$ be the orbit dimension.  A wedge of
$r$ independent infinitesimal action fields is an invariant regular
relative anticanonical section on the generic fibre; commutativity of
$T$ gives invariance.  A rational frame of $-K_Y$ and the canonical
Jacobian section $s_J$ put it in
\[
 p^*(-K_V)=-K_{W/Y}+f^*(-K_Y)+J.
\]
It has no horizontal poles. Regard this section of $p^*(-K_V)$
birationally as a rational section $\rho$ on $V$.  The Jacobian section
is invariant for the induced canonical actions.  The ratio of any two
invariant sections of $J+mp^*(-K_V)$ is a rational invariant function,
hence belongs to the base field.  The sections $s_J(p^*\rho)^m$ supply the
nonzero element at every $m\geq0$, including $m=0$.  Thus the invariant
rank-one condition holds.

If $Y$ is a point, the full wedge of infinitesimal fields on $V$ already
gives a global invariant anticanonical section.  Otherwise the invariant
case of Lemma~\ref{rb:rank-one} and Lemma~\ref{rb:transfer} gives a
bounded semipositive base bundle and a finite-volume donor metric.
The base has no holomorphic forms by dominant pullback, and
Theorem~\ref{rb:allowed} gives a section with permitted base poles.
Its transferred section is invariant, and its only permitted poles on
$V$ lie in $D=0$.  It is therefore regular.
\end{proof}

\subsection{A direct integral argument on logarithmic orbits}

We next establish the two psh functions needed for transfer by an
integral calculation, independently of the invariant direct-image
curvature argument.  The calculation uses the variance estimate of
Brascamp--Lieb \cite{BrascampLieb1976}; we include the needed proof on a
convex ball to make the limiting procedure explicit.
The integral calculation is a complex Pr\'ekopa argument, related to
Berndtsson's formulations \cite{Berndtsson1998,Berndtsson2006}; the
supremum assertion is the corresponding minimum principle of Kiselman
\cite{Kiselman1978}. The tube-domain formulation in
\cite[Theorem~1.1]{DengWangZhangZhou2018} gives a direct comparison for
that second assertion. We prove both statements here because the
almost-everywhere finite integrals and the exhaustion limits are
precisely the data used by the transfer construction.

\begin{lemma}[Integral and supremum weights]\label{g:prekopa}\label{rb:prekopa}
Let $r\geq0$ be an integer, let $U$ be a complex manifold, and let $\Phi(z,x)$ be a smooth function
on $U\times\mathbb R^r$ such that $\Phi(z,\operatorname{Re}w)$ is
plurisubharmonic on $U\times\mathbb C^r$.  Suppose
\[
                     0<\int_{\mathbb R^r}e^{-\Phi(z,x)}\,dx<\infty
\]
for almost every $z$.  Then the almost-everywhere functions
\[
 -\log\int_{\mathbb R^r}e^{-\Phi(z,x)}\,dx,
 \qquad
 -\log\sup_{x\in\mathbb R^r}e^{-\Phi(z,x)}
\]
have psh representatives, whenever the second expression is taken on
the locus where its supremum is finite and positive.  Under the stated
integrability assumption that supremum is finite and positive for
almost every $z$.
\end{lemma}

\begin{proof}
For $r=0$, give $\mathbb R^0$ its unit mass. Both expressions equal $\Phi(z)$, so the assertion is immediate. Assume $r>0$.
Plurisubharmonicity implies convexity of $x\mapsto\Phi(z,x)$.
An integrable positive smooth log-concave function on $\mathbb R^r$
is bounded.  To see this, choose a fixed ball on which its value is at
least $c>0$.  If values $a_j$ at points $x_j$ tended to infinity,
log-concavity on the midpoint ball between $x_j$ and that fixed ball
would give a lower bound $(ca_j)^{1/2}$ on a ball of a fixed positive
volume.  Its integral would tend to infinity, a contradiction.
This proves the last assertion.

Fix a bounded real ball $D$ and first replace $\Phi$ by
$\Phi+\epsilon|x|^2$, so the Hessian in $x$ is positive definite. Until $\epsilon$ is sent to zero,
write $\Phi$ for this regularized weight. Restrict the base variable to a complex disc.  Let $\mathbb E$ denote
expectation for the probability density proportional to $e^{-\Phi}$
on $D$.  Differentiating under the integral gives
\begin{equation}\label{rb:variance}
 \frac{\partial^2}{\partial z\partial\bar z}
          \left(-\log\int_D e^{-\Phi}\,dx\right)
   =\mathbb E\Phi_{z\bar z}
       -\mathbb E|\Phi_z-\mathbb E\Phi_z|^2.
\end{equation}
For a smooth real weight $b$ whose Hessian $b_{xx}$ is positive definite
on $\bar D$, put
\[\mathbb E_b F=\frac{\int_D F e^{-b}dx}{\int_D e^{-b}dx}.\]
The Brascamp--Lieb bound is
\begin{equation}\label{rb:bl}
 \mathbb E_b|v-\mathbb E_bv|^2
       \leq\mathbb E_b\bigl(b_{xx}^{-1}(\nabla v,\nabla\bar v)\bigr).
\end{equation}
Here it is enough to prove it for smooth real mean-zero $v$ and then
add the estimates for real and imaginary parts. The bounded-domain
form and its Neumann identity are also proved in
\cite[Theorems~1.1 and~1.2(1), and Section~3.3]{KolesnikovMilman2017}.
The preceding addition of \(\epsilon|x|^2\) provides the stated positive
definiteness; the limit \(\epsilon\downarrow0\) is taken afterwards.

Solve the weighted Neumann problem
\[
             (-\Delta+\nabla b\cdot\nabla)u=v,
             \qquad \partial_{\mathbf n}u=0.
\]
The mean-zero condition is precisely its solvability condition on the
ball.  Weighted integration by parts and the differentiated equation
give
\[
 \int_Dv^2e^{-b}
  =\int_D|\nabla^2u|^2e^{-b}
     +\int_Db_{xx}(\nabla u,\nabla u)e^{-b}
     +\int_{\partial D}\mathrm{II}(\nabla u,\nabla u)e^{-b}.
\]
The boundary second fundamental form $\mathrm{II}$ is nonnegative
because $D$ is convex, and $\nabla u$ is tangential there.  Also
$\int_Dv^2e^{-b}=\int_D\nabla v\cdot\nabla u\,e^{-b}$.
Cauchy--Schwarz for the quadratic form $b_{xx}$ bounds the latter
integral by
\[
 \left(\int_Db_{xx}^{-1}(\nabla v,\nabla v)e^{-b}\right)^{1/2}
 \left(\int_Db_{xx}(\nabla u,\nabla u)e^{-b}\right)^{1/2}.
\]
Combining the two identities gives \eqref{rb:bl}, including the case
of zero variance by continuity.

Apply \eqref{rb:bl} with $b(x)=\Phi(z,x)$ and $v=\Phi_z$.  The complex Hessian of
$\Phi(z,\operatorname{Re}w)$ is positive semidefinite.  Its Schur
complement gives
\[
         \Phi_{z\bar z}\geq
         \Phi_{zx}(\Phi_{xx})^{-1}\Phi_{x\bar z}.
\]
The factors $1/2$ in complex vertical derivatives cancel against the
factor $1/4$ in the vertical Hessian.  Equations~\eqref{rb:variance}
and \eqref{rb:bl} prove plurisubharmonicity of
$-\log\int_D e^{-\Phi}dx$.  Decreasing $\epsilon$ to zero preserves
this conclusion.  Exhausting $\mathbb R^r$ by balls gives a decreasing
sequence of psh functions; its almost-everywhere finite limit is the
first required representative.

For the supremum, use the high-power integral argument of
\cite[p.~1635, following Theorem~1.2]{Berndtsson2006}: take the
probability-normalized Lebesgue measure on $D$ and replace $\Phi$ by
$l\Phi$, with integers $l\geq1$. The functions
\[
       -\frac1l\log\left(\frac1{|D|}\int_D e^{-l\Phi}dx\right)
\]
are psh and decrease with $l$, by monotonicity of $L^l$ norms.  Their
finite limit is $-\log\sup_D e^{-\Phi}$.  Increasing the balls again
gives a decreasing psh limit, finite almost everywhere by the first
paragraph.  This is the second representative.
\end{proof}

\begin{proof}[Orbit-integral proof of Theorem~\ref{rb:invariant}]
If the torus acts trivially, apply Theorem~\ref{g:allowed} with
$D=0$ and $A=L=-K_V$. Assume the action is nontrivial.
Give $L=-K_V$ the same compact-torus-invariant smooth semipositive
metric and let $\mu$ be its smooth positive volume. Use the maps
$p:W\to V$, $f:W\to Y$ and intermediate model $\bar V$ recalled
in the first proof.  Divide out the
ineffective kernel of the torus action.  The effective torus has
trivial generic stabilizer.  Indeed a subgroup fixing the generic
point would act trivially on the function field and hence on $V$;
subgroups of a split torus are determined by their character lattices,
so the geometric generic stabilizer is such a fixed subgroup.
The generic orbit over the rational quotient field is consequently
a torsor under a split torus.  Hilbert's Theorem~90 trivializes it.
Spreading out gives a dense orbit family $T\times Y^0$ over a nonempty
base open set.  These quotient and triviality statements use
\cite{Rosenlicht1956,BellGhiocaReichstein2017} for the quotient and
\cite[Corollary~3.47 and Proposition~3.50]{MilneAlgebraicGroups2017}
for split-torus torsor triviality.

On this product take logarithmic coordinates
$t_j=e^{x_j+i\theta_j}$ on $T$, with $x\in\mathbb R^r$ and
$\theta\in(\mathbb R/2\pi\mathbb Z)^r$.  Wedge the logarithmic action fields
with a base coordinate anticanonical frame to obtain a frame $\sigma$
of $L$ along the product.  Its weight
\[
                  \Phi(y,x)=-\log|\sigma|^2
\]
is smooth and psh in $(y,x+i\theta)$, and is independent of $\theta$
by compact invariance.  In one coordinate,
\[
 \frac{i}{2}\frac{dt_j}{t_j}\wedge
             \frac{d\bar t_j}{\bar t_j}=dx_j\wedge d\theta_j.
\]
The top form dual to the logarithmic anticanonical frame therefore
identifies the volume with a fixed positive constant times
$e^{-\Phi}\,dx\,d\theta\,d\lambda_y$. Angular integration contributes
a positive constant independent of $y$.
The complement of the orbit family is algebraic and has measure zero.
Consequently the pushforward density of $\mu$ is
\[
                  g_0(y)=C\int_{\mathbb R^r}e^{-\Phi(y,x)}dx
\]
almost everywhere.  It is finite almost everywhere by Fubini.

Choose a nonzero rational base anticanonical frame.  It differs from
the coordinate frame by a base factor and gives an invariant rational
section $\rho$ of $L$.  For almost every general fibre, the supremum of
$|\rho|^2$ on its dense orbit is finite by Lemma~\ref{rb:prekopa}.
It follows that $\rho$ has no horizontal pole on the normal intermediate model $\bar V$:
a pole along a horizontal prime would make its norm unbounded near
that prime on a general compactified fibre, and the dense orbit
approaches its general points.  The smooth metric is comparable to a
positive smooth frame there, so such unboundedness cannot be hidden
in the norm. Normality therefore makes the section regular on the
complete generic fiber of \(\bar V\to Y\), and its pullback to \(W\)
remains regular.

The supremum on the dense orbit equals the essential supremum for
transfer, since the rational section is now regular on the generic
fibre and the metric is continuous.  Lemma~\ref{rb:prekopa} proves
both $-\log g_0$ and $-\log S_\rho$ psh on the quotient open set;
nonzero base frame factors contribute only pluriharmonic terms.
If the quotient is a point, the same no-horizontal-pole argument
already says that $\rho$ is a regular invariant section. Assume now
$\dim Y>0$. Apply Lemma~\ref{g:transfer} with
$D=0$, $A=L=-K_V$, $k=1$, and this rational section $\rho$.
The preceding calculation supplies its two psh weights. Vertical
normalization, properness, and the bounded input metric give the local
upper bound for the normalized section norm; a component attaining
the vertical minimum gives the complementary lower bound. Thus the
transfer produces a bounded metric on its line $B$ and a finite-volume
metric on the different line $-K_Y+D_Y$.
Theorem~\ref{g:allowed} applies on $Y$, which inherits the absence of
positive-degree holomorphic forms by dominant pullback. The output has
the explicit form $(\psi\circ f)(p^*\rho)^m$ on $W$ for a rational
base function $\psi$. Its birational pushdown to $V$ is therefore invariant. Since the initial boundary was
$D=0$, the support conclusion of transfer removes all divisorial poles
on $V$. This completes the orbit-integral proof.
\end{proof}

\subsection{The rationally connected variant}
\label{rb:rc-invariant}
If $V$ is rationally connected, the first proof also closes using only
Proposition~\ref{rb:rc-method} as its nonvanishing input.  Here is the
exact substitution.  For the trivial action, apply that proposition
with $A=L=-K_V$ and $D=0$.  For a positive-dimensional orbit, construct
$\rho$ and the invariant rank-one spaces exactly as above, retaining
$m=0$.  The rational quotient base $Y$ is rationally connected, since
it is a smooth projective model of a dominant rational image of $V$.
The two metric lemmas give $B$ with bounded semipositive weights and
$A_Y=-K_Y+P$ with finite distinguished volume, where $P$ is the
allowed-prime divisor.  Thus Proposition~\ref{rb:rc-method} applies on
$Y$ with target $B$ and boundary $P$.  Its section transfers to an
invariant section on $V$, and the support assertion gives no pole
because the original boundary is zero.  The point quotient is handled
by the full wedge of action fields, as in the first proof.  Consequently
the maximal-fibre-space argument supplies invariant anticanonical
nonvanishing on every rationally connected smooth projective $V$ with
a smooth semipositive anticanonical metric.
